% AUTO-GENERATED FILE. Source: pub/claims.toml
Evaluation remains the least consolidated part of explainable artificial intelligence (XAI). While model-agnostic post-hoc methods such as LIME, SHAP, Anchors, and DiCE are widely used, the criteria used to judge them vary sharply across papers, toolkits, and application domains. Some studies emphasize faithfulness proxies, others prioritize stability, sparsity, or human plausibility, and newer work introduces large language models (LLMs) as semantic evaluators without a unified framework for locating those judgments within the broader evaluation landscape. This paper presents a structured survey and taxonomy of evaluation metrics for model-agnostic explainability, focusing on how technical proxy metrics, benchmark-grounded metrics, human evaluation, and LLM-based semantic evaluation relate to one another. Using a cleaned 24-study coded review corpus derived from the thesis literature matrix, we organize XAI evaluation along four axes: evaluation target, evidence source, quality property, and task context. The resulting taxonomy clarifies where current metrics overlap, where they leave construct gaps, and why semantic evaluation should be treated as a complementary layer rather than a replacement for robustness- and fidelity-oriented measurement. The paper closes by mapping the taxonomy back to the thesis framework, motivating the benchmark design of Paper B and the evaluator-validation agenda of Paper A.
